% Rev 8 -- Low-Level Embedded Firmware and Engineering Data Analysis
\documentclass[10pt,letterpaper]{article}

\usepackage[margin=0.50in]{geometry}
\usepackage[T1]{fontenc}
\usepackage[utf8]{inputenc}
\usepackage[default,type1]{sourcesanspro}
\usepackage[final]{microtype}
\usepackage[explicit]{titlesec}
\usepackage{enumitem}
\usepackage{tabularx}
\usepackage{array}
\usepackage{ragged2e}
\usepackage{iftex}
\usepackage[hidelinks]{hyperref}

\ifPDFTeX
  \input{glyphtounicode}
  \pdfgentounicode=1
\fi

\hypersetup{
  pdftitle={Connor Savugot Resume},
  pdfauthor={Connor Savugot},
  pdfsubject={Low-Level Embedded Firmware and Engineering Data Analysis},
  pdfkeywords={embedded software, embedded C, C++, MATLAB, Python, real-time systems, RTOS, FreeRTOS, bare metal, power control, ADC filtering, ADC scaling, MCP3428, fixed-point arithmetic, CAN data analysis, DSP, dsPIC33CK, PIC32, TMS320, JTAG, hardware software integration, multiprocessor systems, SoC, Arm, Embedded Linux, firmware validation, Git}
}

\pagestyle{empty}
\setcounter{secnumdepth}{0}
\setlength{\parindent}{0pt}
\setlength{\parskip}{0pt}
\hfuzz=0.4pt

\titleformat{\section}
  {\fontsize{11.3}{11.9}\selectfont\bfseries\uppercase}
  {}{0pt}{#1\hspace{0.12cm}\titlerule[0.65pt]}
\titlespacing{\section}{0pt}{0.14cm}{0.05cm}

\newlist{resumebullets}{itemize}{1}
\setlist[resumebullets]{
  label=\textbullet,
  leftmargin=0.42cm,
  topsep=0.004cm,
  itemsep=0.004cm,
  parsep=0pt,
  partopsep=0pt
}

\newcommand{\companyheader}[1]{%
  {\fontsize{10.4}{10.8}\selectfont\bfseries #1}\par\vspace{-0.03cm}
}

\newcommand{\resumeentry}[2]{%
  \begin{tabularx}{\textwidth}{@{}X>{\raggedleft\arraybackslash}p{2.40in}@{}}
    \textbf{#1} & #2
  \end{tabularx}\vspace{-0.075cm}
}

\newcommand{\roleentry}[2]{%
  \begin{tabularx}{\textwidth}{@{}X>{\raggedleft\arraybackslash}p{2.40in}@{}}
    \textbf{#1} & #2
  \end{tabularx}\vspace{-0.075cm}
}

\newcommand{\descriptorentry}[1]{%
  \begin{tabularx}{\textwidth}{@{}X>{\raggedleft\arraybackslash}p{2.40in}@{}}
    \textit{#1} &
  \end{tabularx}\vspace{-0.075cm}
}

\begin{document}
\fontsize{9.65}{10.50}\selectfont
\setlength{\RaggedRightRightskip}{0pt plus 5em}
\RaggedRight
\emergencystretch=1em
\hyphenpenalty=10000
\exhyphenpenalty=10000

\begin{center}
  {\fontsize{21.5}{22.2}\selectfont\bfseries Connor Savugot}\\[1.5pt]
  {\bfseries B.S. in Computer Engineering | North Carolina State University | Graduated May 2026}\\[1.5pt]
  {\fontsize{9.6}{10.2}\selectfont
    \href{mailto:consav04@gmail.com}{consav04@gmail.com}
    \enspace|\enspace \href{tel:+18285419487}{+1-828-541-9487}
    \enspace|\enspace Murphy, NC
    \enspace|\enspace \href{https://linkedin.com/in/connorsavugot}{linkedin.com/in/connorsavugot}
    \enspace|\enspace \href{https://github.com/Clem085}{github.com/Clem085}}
\end{center}
\vspace{-0.18cm}

\section{Technical Profile}
Embedded firmware engineer developing real-time C for power, control, and communications systems across 16- and 32-bit MCUs, FreeRTOS, and bare-metal DSPs. Builds MATLAB/C++/Python engineering tools, brings up peripherals from hardware documentation, and collaborates with a distributed firmware team on hardware and embedded Linux integration.

\section{Technical Skills}
\textbf{Languages:} Embedded C, C++, MATLAB, Python, Bash\\
\textbf{Embedded \& Real-Time:} Bare-metal firmware, FreeRTOS, POSIX threads, interrupts/timers, peripheral/HAL bring-up, boot/update/recovery systems\\
\textbf{Processors:} dsPIC33CK (16-bit DSC), PIC32, TMS320 DSP, STM32G474, MSP430FR2355; TI AM62 Arm-based SoC\\
\textbf{Analysis, Debug \& Build:} ADC scaling/filter analysis, MATLAB data analysis, CAN-log parsing, JTAG, register inspection, Code Composer Studio, GCC/Make, Git/GitLab, Docker, oscilloscopes, logic/protocol analyzers\\
\textbf{Systems \& Interfaces:} Embedded Linux, Yocto/Arago, device trees, systemd; PWM, ADC, I2C, SPI, UART/RS-232; CAN 2.0, CAN FD, J1939, DroneCAN; IPMI/IPMB

\section{Engineering Experience}
\companyheader{AEGIS Power Systems}
\roleentry{Embedded Firmware Engineer}{May 2026--Present}
\begin{resumebullets}
  \item Develop low-level C firmware for 16-bit dsPIC33CK DSCs and bring up peripherals one subsystem at a time from datasheets, reference manuals, and schematics; isolate register, timing, configuration, and electrical faults using JTAG, targeted tests, oscilloscopes, and logic analyzers.
  \item Parsed captured CAN output in MATLAB to determine per-channel ADC scaling and compare several filter functions by steady-input variation and time to follow large electrical changes; selected rolling-average and exponential moving average (EMA) approaches and implemented the rolling average in embedded C using a circular buffer and rolling sum.
  \item Implemented and validated VITA 46.11 Tier 2 management firmware for a VPX power supply, enabling the chassis manager to read IPMI 2.0 sensor, FRU/inventory, and event data over an I2C-based IPMB link.
  \item Integrate battery hardware and CAN-connected power devices with a TI AM62 Arm-based embedded Linux controller; build Yocto/Arago images and kernels, update device trees, develop external-watchdog services, and trace faults across power, wiring, CAN, boot, configuration, and code.
\end{resumebullets}

\roleentry{Embedded Systems Intern}{May--Aug 2025}
\roleentry{Computer Science Intern}{May--Aug 2024}
\descriptorentry{Selected work across both internships}
\begin{resumebullets}
  \item Developed bare-metal C power-control firmware on a TMS320 DSP for transformer-coupled bidirectional buck--boost supplies using 12 PWM A/B pairs (24 physical outputs); configured each A channel's period/frequency, duty cycle, and relative phase while each B output followed or inverted its paired A signal, then validated ADC/PWM behavior in Code Composer Studio.
  \item Developed embedded C firmware across a two-processor power system: a PIC32 running FreeRTOS coordinated I2C, SPI, CAN/J1939, UART, and update behavior while the bare-metal TMS320 DSP executed real-time converter control.
  \item Built a resilient update/recovery workflow: a C++ CLI accepted HEX files, a Python TCP sender handled chunking and checksums, and the PIC32 staged new firmware and a golden image in external SPI flash, programmed and verified the DSP over UART, and controlled boot.
  \item In 2024, built a restricted SSH/chroot interface and Python/Bash commands for a resource-constrained TI Yocto/Arago power system; reverse-engineered RS-232 commands/checksums for programmable loads and relay controllers and helped assemble the associated test circuit.
\end{resumebullets}

\section{Selected Engineering Projects}
\resumeentry{EV Active Sensor Adapter | Battery Monitoring \& CAN}{NC State University}
\begin{resumebullets}
  \item Developed STM32G474/BQ79616 bring-up and integration for a team-built EV sensor adapter spanning cell-voltage monitoring, ten thermistors sampled by the STM32 ADC, UART, extended-ID CAN telemetry, subsystem-level fault isolation, and custom-PCB hardware.
  \item Implemented BQ79616 wake/configuration, protocol framing, CRC validation, cell reads, and fault diagnostics; coordinated requirements, staged validation, risk tracking, design reviews, and handoff documentation across the team.
\end{resumebullets}

\resumeentry{Real-Time Systems \& Formal Verification | C, Linux, Lustre}{Spring 2026}
\begin{resumebullets}
  \item Implemented periodic POSIX-thread tasks in C using CPU affinity, absolute-time releases, priority-inheritance mutexes, and timestamp traces; analyzed schedulability, plotted execution timelines, and used bounded model checking on Lustre models to correct temporal monitors and check invariants within configured bounds.
\end{resumebullets}

\resumeentry{MATLAB Engineering Analysis \& Machine Learning}{Selected Coursework}
\begin{resumebullets}
  \item Analyzed measured MOSFET CSV data in MATLAB using numerical gradients and polynomial fits to estimate transconductance, threshold voltage, mobility, and channel-length modulation and compare measured and modeled I--V curves.
  \item Implemented MATLAB image-learning workflows: trained a CNN to 99.24\% validation accuracy, adapted SqueezeNet through transfer learning on a five-class coursework dataset, and applied PCA/KLT to visualize eigenfaces and reduce image dimensionality.
\end{resumebullets}

\section{Technical Leadership}
\resumeentry{ECE 306 Embedded Systems | Teaching Assistant \& Course Mentor}{Spring 2025--Spring 2026}
\begin{resumebullets}
  \item Delivered supplemental lectures and debugged embedded C/MSP430FR2355 firmware and custom-PCB projects across three semesters, diagnosing register, timing, peripheral, toolchain, wiring, and soldering failures across diverse student implementations.
\end{resumebullets}

\end{document}
